\documentclass[10pt,twocolumn]{article}
\usepackage[margin=0.85in]{geometry}
\usepackage{amsmath}
\usepackage{amssymb}
\usepackage{booktabs}
\usepackage{hyperref}
\usepackage{listings}
\usepackage{xcolor}

\hypersetup{colorlinks=true, linkcolor=violet, urlcolor=teal, citecolor=violet}

\lstset{
  basicstyle=\ttfamily\footnotesize,
  breaklines=true,
  frame=single,
  columns=fullflexible,
  keepspaces=true,
}

\title{\textbf{PayLane: A Reserve-Gated Stablecoin Treasury on Solana}}
\author{Dipankar Sarkar}
\date{2026}

\begin{document}
\maketitle

\begin{abstract}
Fiat-backed stablecoins are the settlement layer of on-chain payments, yet their
solvency is enforced off-chain: an issuer holds reserves at a bank and promises
that circulating supply never exceeds them. When that promise breaks, holders
learn only after the fact. PayLane moves the promise into the program. It is a
Solana treasury, written in Anchor, whose \texttt{mint} and \texttt{settle}
instructions run a reserve gate that provably cannot push circulating supply past
the reserves attested on-chain. The gate is a small, total function over unsigned
64-bit integers with no floating point and no administrative bypass; a mint of one
unit beyond the attested reserve is rejected deterministically and leaves state
untouched. Circulating supply is never a number the program stores and hopes is
right: it is read from the bound SPL \texttt{Mint.supply} on every instruction, and
redemption burns the holder's tokens, so the ledger cannot drift from the tokens
that actually exist. An attestation that reports less backing than is circulating
is recorded rather than refused, and pauses minting until solvency returns. We port
the ProofBackedUSD reserve-gate idea from an EVM contract to Anchor, factor the
invariant into a pure-Rust module that is unit-tested without the Solana toolchain,
and exercise the instructions themselves against the real SPL Token program under
\texttt{solana-program-test}. The result is a payment rail whose solvency is a
property of the code, checkable by anyone, rather than a claim of the operator.
\end{abstract}

\section{Problem}
A fiat-backed stablecoin issuer maintains two numbers: the reserves it holds and
the supply it has minted. The system is solvent exactly when supply does not
exceed reserves. On every major stablecoin today this inequality is maintained by
policy and audited periodically, not enforced at the point of minting. The mint
authority is a key; whoever holds it can, in principle, mint beyond reserves,
whether by error, by compromise, or by choice. Redemption at par then depends on
trust that this has not happened between audits.

This is the wrong place for the check. Minting is an on-chain action; the reserve
figure can be attested on-chain; and the comparison between them is a single
integer inequality. There is no technical reason the chain cannot refuse an
over-mint the way it refuses a double-spend. The obstacle is that the check has
historically lived in the issuer's back office rather than in the program that
controls the mint. PayLane relocates it.

\section{Related Work}
Collateralized stablecoins such as MakerDAO's DAI enforce over-collateralization
on-chain through per-vault liquidation, a different model in which each borrower's
position is independently backed by volatile crypto collateral~\cite{makerdao}.
PayLane instead targets the fiat-backed model, where a single treasury attests a
scalar reserve figure and the invariant is a global supply ceiling. Transparency
efforts for fiat-backed issuers, such as Circle's published reserve
attestations~\cite{circle-transparency}, make the reserve figure visible but do
not bind minting to it; the mint remains a trusted key. Proof-of-reserves work in
the exchange setting~\cite{proofofreserves} likewise proves a snapshot without
gating future issuance against it. On Solana specifically, the SPL Token program
provides a mint authority abstraction~\cite{spl-token} but delegates all policy to
that authority; Anchor~\cite{anchor} is the framework we use to wrap the mint
authority in a program-derived treasury that interposes the reserve gate. PayLane's
contribution is to make the attested reserve a hard, program-enforced ceiling on
issuance rather than a disclosed number.

\section{Design}
The treasury is a single program-derived account holding four fields relevant to
solvency: the \texttt{authority} that may mint and settle, the \texttt{attestor}
that may update reserves, the bound SPL mint, one unsigned 64-bit
\texttt{attested\_reserve} denominated in the token's smallest unit, and a pause
flag. Circulating supply is deliberately \emph{not} stored: it is projected from
the bound mint's \texttt{supply} field at instruction time. The solvency invariant
is
\[
\texttt{circulating} \le \texttt{attested\_reserve}
\]
and every state transition preserves it.

\paragraph{The gate.} Minting an amount $a$ is allowed only when
$\texttt{circulating} + a \le \texttt{attested\_reserve}$, computed with a checked
addition so that a would-be $2^{64}$ overflow is rejected rather than wrapped.
Settlement lowers circulating supply and cannot break the invariant, but is guarded
against underflow. A new attestation is always recorded, including one that drops
reserves below what is already circulating: that figure is the truth about a
shortfall, and refusing it would leave a stale, higher number on the books while
the issuer kept minting against reserves that no longer exist. Such an attestation
instead pauses the treasury, and every mint is refused until an attestation covers
supply again or holders redeem back under the line. The whole gate is a total
function: every input maps to either a new valid state or a typed rejection, and a
rejection is a no-op on state.

\paragraph{One source of truth.} The gate lives in a pure-Rust module,
\texttt{reserve.rs}, that has no dependency on the Anchor runtime. The on-chain
program imports it verbatim and calls it inside each instruction handler; a
sibling crate imports the same file and unit-tests it with plain \texttt{cargo test}.
The security-critical logic that runs on-chain is therefore byte-identical to the
logic under test, and the tests need neither a validator nor the Solana BPF
toolchain to run.

\paragraph{Separated powers.} The authority and the attestor are distinct keys.
The party that mints cannot raise its own reserve ceiling, and the party that
attests reserves cannot mint. This mirrors the operational separation between a
treasury desk and a proof-of-reserve oracle.

\paragraph{Agent and adapter seam.} An off-chain agent, built on a small reusable
agent core, drives the treasury. It reads treasury state, proposes mints and
settlements, and runs the \emph{same} reserve gate off-chain before submitting a
transaction, so it never broadcasts an instruction the program would reject. Each
action is journaled with a guardrail decision, giving an auditable activity log.
The agent talks to the chain through an adapter seam whose rule is that submission
comes first: nothing is recorded locally until a transaction confirms, and a failed
submission leaves the mirror untouched. This build ships no live client, because
nothing is deployed; when the seam is configured for a cluster and no client is
supplied it refuses to construct rather than run the simulation and call it live.
The deterministic offline mirror is labelled as such, down to the event digests,
so the demo is keyless without claiming to be on-chain.

\section{Hero}
The hero action is an attempted over-mint. With one million dollars attested and
eight hundred thousand circulating, the agent proposes minting three hundred
thousand more. That would put circulating supply at $1{,}100{,}000$ against
$1{,}000{,}000$ of backing. The gate computes $800{,}000 + 300{,}000 = 1{,}100{,}000 >
1{,}000{,}000$ and returns \texttt{ExceedsReserves}. No tokens are minted, no
transaction is produced, and the treasury state is exactly what it was before. The
operator then settles one hundred thousand out of circulation, which frees
headroom, and the identical three-hundred-thousand mint now fits under the ceiling
and is allowed. The point is visceral: solvency is not something the operator
asserts between audits, it is something the program refuses to violate on every
single mint.

\section{Evaluation}
We evaluate the gate directly, since it is the entire security claim, and then the
program that wraps it. Nineteen pure-Rust unit tests exercise the invariant: mints
within reserves, a mint exactly to the ceiling, a mint one unit past it, that a
rejected mint does not mutate state, zero-amount rejection, overflow near $2^{64}$,
settlement lowering supply and freeing headroom, settlement underflow, an
attestation below circulating supply being recorded and pausing the treasury, that
no mint is possible while paused, that a covering attestation or a large enough
redemption unpauses while a partial one does not, the backing ratio reporting the
shortfall honestly and saturating instead of returning a sentinel, and a multi-step
operator sequence that preserves the invariant at every step. They run in the
Anchor crate and, byte for byte, in a sibling crate that needs no Solana toolchain.

Thirteen further tests drive the actual instructions against the real SPL Token
program inside \texttt{solana-program-test}, an in-process bank. These assert on
the SPL \texttt{Mint} account rather than on anything PayLane wrote about itself:
that \texttt{initialize} moves mint authority to the treasury PDA and that a
stranger cannot claim another issuer's mint, that a successful mint moves real
supply and credits the holder while a rejected one moves nothing, that redemption
burns the holder's tokens so settlement cannot manufacture unbacked headroom, that
a recorded reserve loss pauses minting on-chain, and that the attestor and the
minting authority cannot act for each other. A stubbed mint helper fails every one
of them.

A Python suite of forty-four tests re-checks the same invariant in the off-chain
port case for case, pins the Anchor wire format (instruction discriminators and the
\texttt{Treasury} account layout), proves the adapter seam commits nothing when a
submission fails, and exercises the CLI workflow and the guardrailed audit trail
end-to-end.

\begin{table}[h]
\centering
\begin{tabular}{lc}
\toprule
\textbf{Suite} & \textbf{Result} \\
\midrule
Rust reserve invariant (\texttt{cargo test}) & 39 / 39 \\
Rust program vs. SPL Token (\texttt{cargo test}) & 13 / 13 \\
Python agent + seam (\texttt{pytest}) & 44 / 44 \\
\bottomrule
\end{tabular}
\caption{Test results. The Rust suite proves the on-chain gate; the Python suite
proves the off-chain agent agrees with it.}
\end{table}

\section{Limitations}
The reserve gate is the trust-minimized core, but it does not by itself make the
whole system trustless. The attested reserve is only as good as the attestor: a
dishonest or compromised attestor can post a reserve figure not backed by real
custody, and the on-chain gate will honor it. PayLane narrows the trusted surface
from ``trust that supply never exceeds reserves'' to ``trust that the attested
reserve reflects custody,'' which is a smaller and more auditable claim, but it
does not eliminate it; a production deployment would pair the gate with a
proof-of-reserve oracle or a multi-party attestor. Second, and most importantly for
a reader judging what has been demonstrated: the program has \emph{not} been
deployed. Compiling it for deployment needs the Solana toolchain
(\texttt{anchor build} / \texttt{cargo build-sbf}), which is not installed in this
environment, so there is no devnet program id and no explorer-verifiable
transaction. What is demonstrated is the program's behaviour under the real SPL
Token program in an in-process bank, which is the same code path minus the SBF
compilation and the network. Only devnet deployment is in scope; mainnet is not.
Third, the treasury is bound to a single mint and a single attestor key; multi
-currency treasuries and threshold attestation are future work.

\begin{thebibliography}{9}
\bibitem{makerdao} MakerDAO. \emph{The Maker Protocol White Paper}. 2020. \url{https://makerdao.com/en/whitepaper}
\bibitem{circle-transparency} Circle. \emph{USDC Transparency and Reserves}. \url{https://www.circle.com/en/transparency}
\bibitem{proofofreserves} G. Maxwell et al. \emph{Provisions: Privacy-preserving proofs of solvency for Bitcoin exchanges}. ACM CCS, 2015. \url{https://eprint.iacr.org/2015/1008}
\bibitem{spl-token} Solana Labs. \emph{SPL Token Program}. \url{https://spl.solana.com/token}
\bibitem{anchor} Coral / Anchor contributors. \emph{The Anchor Framework}. \url{https://www.anchor-lang.com/}
\end{thebibliography}

\end{document}
